\documentclass[aspectratio=169]{beamer}
\beamertemplatenavigationsymbolsempty
\usetheme{Madrid}
\usepackage{graphicx, hyperref, amsmath, animate, tabularx, bm}
\usecolortheme{beaver}
\usefonttheme{professionalfonts}
\graphicspath{{/Users/soumen/Desktop/IITB/ms_research/MS_Report/BPS1/}}

\usepackage[backend=bibtex, style=verbose, sorting=none, autocite=footnote]{biblatex}
\addbibresource{References.bib}
\renewcommand*{\bibfont}{\scriptsize}
\setbeamerfont{footnote}{size=\tiny}

\title[MS by Research: BPS (Autumn 2024)]{Improving Downstream Task Performance in Multilingual LLMs by Intervening Language Specific Neurons}
\author[Soumen Mondal]{\textcolor{blue}{Soumen Kumar Mondal} \\\texttt{23m2157@iitb.ac.in}}
\institute[IIT Bombay]{Guide: Prof.~Preethi Jyothi\\\;\\MS by Research in Data Science and Artificial Intelligence\\Center for Machine Intelligence and Data Science (CMInDS)\\Indian Institute of Technology Bombay}
\date{December 26, 2024}

\begin{document}
	% Title page frame
	\begin{frame}
		\titlepage
	\end{frame}
	
\begin{frame}{Introduction}
	\begin{block}{\scriptsize Problem Statement}\scriptsize
		\begin{itemize}
			\item \textbf{Objective:} 
			\begin{itemize}\scriptsize
				\item Investigate the representation of language-specific information in large language models (LLMs).
				\item Develop methods to identify and utilize language-specific neurons for multilingual tasks.
				\item Enhance performance for low-resource languages by employing techniques such as activation probability entropy and LoRA-based fine-tuning.
			\end{itemize}				
			\item \textbf{Importance:}
			\begin{itemize}\scriptsize
				\item Ensure fair and inclusive NLP systems by bridging the performance gap between high-resource and low-resource languages.
				\item Advance understanding of the internal mechanisms of multilingual LLMs.
				\item Contribute to the development of efficient methods for improving multilingual performance.
			\end{itemize}
		\end{itemize}
	\end{block}
	\begin{block}{\scriptsize Presentation Outline}\scriptsize
		\begin{itemize}
			\item Review of existing methods for identifying and utilizing language-specific neurons.
			\item Experimental design: fine-tuning configurations and intervention strategies.
			\item Results and analysis for multilingual tasks (e.g., XNLI task).
			\item Conclusion and future work directions.
		\end{itemize}
	\end{block}
\end{frame}
	
	\begin{frame}{Methodology: Definition of Neuron}
		\begin{columns}
			\column{0.5\textwidth}
			\begin{block}{\scriptsize What is a Neuron?}\scriptsize
				\begin{itemize}
					\item A \textbf{neuron} corresponds to a specific dimension of the activation vector $\mathbf{h}^i$, represented as:
					\begin{equation}
						\mathbf{h}^i = [ \text{act\_fn}(\hat{\mathbf{H}}^i \mathbf{W}_1^i) \otimes (\hat{\mathbf{H}}^i \mathbf{W}_3^i) ] \cdot \mathbf{W}_2^i,
					\end{equation}
					\item The average activation (against \citeauthor{wikidump} sentence) of a neuron $(i, j)$ across token positions is computed as:
					\begin{equation}
						\bar{h}_{i,j} = \frac{1}{T} \sum_{t=1}^{T} h_{i,j}^t.
					\end{equation}
					\item Language-specific activations are computed as:
					\begin{equation}
						h_{i,j}^l = \frac{1}{|S_l|} \sum_{s \in S_l} \bar{h}_{i,j}^{(s)},
					\end{equation}
					where $S_l$ is the set of all sequences for language $l$.
				\end{itemize}		  
			\end{block}
			\column{0.4\textwidth}
			\begin{block}{\scriptsize MLP Architecture of Llama3 [\citeauthor{grattafiori2024llama} et al.]}
				\begin{figure}
					\centering
					\includegraphics[width=\textwidth]{Chapters/Chapter_3/mlp.pdf}
				\end{figure}
			\end{block}
		\end{columns}
	\end{frame}
	
	
\begin{frame}{Methodology: LAPE}
	\begin{columns}
	\column{0.55\textwidth}
	\begin{block}{\scriptsize LAPE Neurons [\citeauthor{tang2024languagespecific}]}\scriptsize
		\begin{itemize}
			\item To identify neurons with significant language-specific behavior, a threshold (for Wikipedia dataset) is defined as:
			\begin{equation}
				\tau = \text{Quantile}_{0.95}\left(p_{i,j}^l \, | \, \forall i, j, l\right).
			\end{equation}
			\item A neuron $(i, j)$ is considered as language neuron if:
			\begin{equation}
				\exists l \in \{1, \dots, k\} \, \text{such that} \, p_{i,j}^l > \tau
			\end{equation}
			
			\item Only 1\% ($m = 0.01 \times 4d \times L$) of neurons with lowest LAPE values are selected for as language neurons.
			\begin{equation}
				LNS = \text{argsort}_m \left\{\text{LAPE}(i,j) \,|\, \forall i, j, \exists l \,\, s.t. \,\,  p_{i,j}^l > \tau \right\}
			\end{equation}
			\begin{equation}
				NTL(i, j) = \left\{ l \, | \, p_{i,j}^l > \tau \land (i,j) \in LNS \right\}
			\end{equation}
		\end{itemize}		  
	\end{block}
	\column{0.4\textwidth}
	\begin{block}{\scriptsize LAPE Workflow}
		\begin{figure}
			\centering
			\includegraphics[width=\textwidth]{Chapters/Chapter_3/lape.pdf}
		\end{figure}
	\end{block}
\end{columns}
	\end{frame}
	
\begin{frame}{Methodology: Activation Based Methods}
	\begin{columns}
		\column{0.5\textwidth}
		\begin{block}{\scriptsize Neuron Relevance}\scriptsize
			\begin{itemize}
				\item \textbf{Activation Probability 90p:} Measures the probability of a neuron's activation exceeding the 90th percentile value for its layer.
				\item Relevance for a neuron \((i, j)\) in language \(l\) is computed from its activation against Wikipedia dataset:
				\begin{equation}
					P_q(h_{i,j}^l) = \inf \left\{x \in \mathbb{R} : CDF_{h_{i,j}^l}(x) \geq q \right\}
				\end{equation}
				\begin{equation}
					r_{i,j}^l = \mathbb{E}\left[\mathbb{I}\left(h_{i,j}^l > P_{0.9}(h_{i,j}^l)\right)\right]
				\end{equation}
				\item Aggregates relevance values across neurons:
				\begin{equation}
					R^l = \{r_{i,j}^l \, | \, i \in \{1, \dots, L\}, \, j \in \{1, \dots, 4d\}\}.
				\end{equation}
			\end{itemize}		  
		\end{block}
		\column{0.4\textwidth}
		\begin{block}{\scriptsize Activation Based Neurons}\scriptsize
			\begin{itemize}
				\item \textbf{Selection Process:} Identify top neurons with the highest relevance for each language.
				\item Select top $0.125\%$ neurons per language based on their maximum relevance:
				\begin{equation}
					 m = \frac{0.125 \times L \times 4d}{100}
				\end{equation}
				\begin{equation}
					LTN(l) = \text{argsort\_desc}_m \left\{r_{i,j}^l \, | \, \forall i, j\right\}
				\end{equation}
				\item This process ensures that neurons with the highest language-specific relevance are identified for each language.
			\end{itemize}
		\end{block}
	\end{columns}
\end{frame}

	
\begin{frame}{Methodology: Neuron Fine-Tuning using LoRA}
	\begin{columns}
		\column{0.55\textwidth}
		\begin{block}{\scriptsize LoRA Fine-Tuning for a Task}\scriptsize
			\begin{itemize}
				\item \textbf{LoRA Mechanism:} Fine-tunes a subset of parameters in MLP by low-rank decomposition [\citeauthor{lora}]:
				\begin{equation}
					[\text{masked\_A}]_{ij} = 1, \forall i \in \{1, \dots, d\}, \, \forall j \in \{1, \dots, r\}
				\end{equation}
				\begin{equation}
					[\text{masked\_B}]_{ij} = 
					\begin{cases} 
						0, & \text{if } j \in \text{Frozen Neurons}, \\
						1, & \text{otherwise}.
					\end{cases}
				\end{equation}
				\item To freeze specific neurons, columns of B corresponding to the frozen neurons are masked with zeros. This prevents any updates to these neurons during fine-tuning.
				\begin{equation}
					\Delta W_{i,j} = 0, \quad \forall i \in {1, \dots, d}, \forall j \in \text{Frozen Neurons}
				\end{equation}
				\item LoRA fine-tunes the MLP layer for the trainable language neurons and attention layer (\(q, k, v, o\) projections). The task-specific classifier head is fully trained.

			\end{itemize}		  
		\end{block}
		\column{0.4\textwidth}
		\begin{block}{\scriptsize Targeted LoRA Fine-Tuning Example}
			\begin{figure}
				\centering
				\includegraphics[width=\textwidth]{Chapters/Chapter_3/lora.pdf}
			\end{figure}
		\end{block}
	\end{columns}
\end{frame}

\begin{frame}{Experiments: Dataset and Baseline}
	\begin{columns}
		\column{0.45\textwidth}
		\begin{block}{\scriptsize Task and Dataset: XNLI}\scriptsize
			\begin{itemize}
				\item \textbf{XNLI Task:} Multilingual natural language inference task to classify relationships between premise and hypothesis into \textit{entailment}, \textit{neutral}, or \textit{contradiction} [\citeauthor{xnli}].
				\item \textbf{Dataset:} Sentence pairs annotated across 15 languages. Focused on English (\textit{en}) as source language and Vietnamese (\textit{vi}), Hindi (\textit{hi}), and Urdu (\textit{ur}) as target languages.
				\item For the experiments, a fraction \(\mathcal{F}\) is used to train the model as $|\mathcal{D}_l^\text{train}| = 400K$.
				\begin{equation}
					|\mathcal{F}| = \text{frac} \times |\mathcal{D}_l^\text{train}|
				\end{equation}
				where \(\text{frac} = 0.25\) is the fraction of the training dataset used.
				\item \textbf{Objective:} Improving the zero-shot performance for low-resource languages (\textit{vi}, \textit{hi}, \textit{ur}).
			\end{itemize}		  
		\end{block}
		\column{0.45\textwidth}
		\begin{block}{\scriptsize Baseline Experiments (No Eval Time Interventions!)}\scriptsize
			\begin{itemize}
				\item \textbf{Zero-Shot Baseline:} Fine-tune on English (\textit{en}) XNLI dataset. Evaluate zero-shot transfer to target languages (\textit{vi}, \textit{hi}, \textit{ur}):
				\begin{equation}
					\text{Accuracy} = \frac{\sum_{i=1}^N \mathbb{I}(\hat{y}^i = y^i)}{N}
				\end{equation}
				where \(N\) is the total test examples.
				\item \textbf{Direct Fine-Tuning Baseline:} Fine-tune on target language such as Vietnamese (\textit{vi}) XNLI dataset to establish the upper bound for \textit{vi} performance. Accuracy remains the evaluation metric.
			\end{itemize}
		\end{block}
	\end{columns}
\end{frame}

\begin{frame}{Experiments: Goals and Fine-Tuning Strategy}
	\begin{columns}
		\column{0.55\textwidth}
		\begin{block}{\scriptsize Experiment Goals}\scriptsize
			\begin{itemize}
				\item \textbf{Objective:} Improve zero-shot performance on target language (\textit{vi}) without using target language training data.
				\item \textbf{Constraints:}
				\begin{itemize}\scriptsize
					\item No target language training data (\textit{vi}) allowed.
					\item Only English (\textit{en}) XNLI training data available.
					\item Use model surgery techniques (e.g., neuron fine-tuning, activation pattern interventions).
				\end{itemize}
				\item \textbf{Target Metric:}
				\begin{equation}
					\Delta_\text{ZS} = \text{Accuracy}(\mathcal{M}_\text{experiment}) - \text{Accuracy}(\mathcal{M}_\text{en})
				\end{equation}
				\item Maximize \(\Delta_\text{ZS}\) to close the performance gap between zero-shot and direct fine-tuning.
				
			\end{itemize}		  
		\end{block}
		\column{0.4\textwidth}
		\begin{block}{\scriptsize Fine-Tuning Setups}\scriptsize
			\begin{itemize}
				\item \textbf{Setup A:} Standard fine-tuning using LoRA applied to \(q, k, v, o\) projection matrices and the classification head. Evaluated with statistical activation substitutions (\(\mu, 90p, 95p\), etc.).
				\item \textbf{Setup B:} Fine-tuning with LoRA targeting English-specific neurons (\(\mathcal{N}_\text{en}\)) in the MLP block.
				\item \textbf{Setup C:} Fine-tuning with LoRA targeting Vietnamese-specific neurons (\(\mathcal{N}_\text{vi}\)) in the MLP block.
				\item \textbf{Setup D:} Combined fine-tuning targeting both English (\(\mathcal{N}_\text{en}\)) and Vietnamese (\(\mathcal{N}_\text{vi}\)) neurons with LoRA.
				
			\end{itemize}
		\end{block}
	\end{columns}
\end{frame}

\begin{frame}{Experiments: Neuron Distribution and Intervention}
	\begin{columns}
		\column{0.55\textwidth}
		\begin{block}{\scriptsize Language Neuron Distribution}\scriptsize
			\begin{itemize}
				\item The number of language-specific neurons for a language $l$ is computed as:
				\begin{equation}
					N^l = \sum_{i=1}^{L} \sum_{j=1}^{4d} \mathbb{I}\left(p_{i,j}^l > \tau \land (i,j) \in LNS \right)
				\end{equation}
				\item Neuron overlap between two languages $l_1$ and $l_2$ is given by:
				\begin{equation}
					O_{l_1, l_2} = \sum_{i=1}^{L} \sum_{j=1}^{4d} \mathbb{I}\left(p_{i,j}^{l_1} > \tau \land p_{i,j}^{l_2} > \tau \land (i,j) \in LNS \right)
				\end{equation}
				\item Layer-wise neuron distribution for a language $l$ is calculated:
				\begin{equation}
					L_l(i) = \sum_{j=1}^{4d} \mathbb{I}\left(p_{i,j}^l > \tau \land (i,j) \in LNS \right)
				\end{equation}
			\end{itemize}		  
		\end{block}
		\column{0.4\textwidth}
		\begin{block}{\scriptsize Perplexity Change}\scriptsize
			\begin{itemize}
				\item Perplexity measures the uncertainty of a model's predictions, defined as:
				\begin{equation}
					\text{PPX}(\mathbf{x}) = \exp\left(-\frac{1}{T} \sum_{t=1}^T \log P(x_t \,|\, x_{<t})\right)
				\end{equation}
				\item Perplexity change ($\text{PPXC}(i,j)$) quantifies the impact of interventions at language $i$ on language $j$:
				\begin{equation}
				\begin{aligned}
					\text{PPXC}(i,j) &= \text{PPX}(j \,|\, \text{Intervention at } i) \\
					&- \text{PPX}(j \,|\, \text{No Intervention at } i)
				\end{aligned}
			\end{equation}
			\end{itemize}
		\end{block}
	\end{columns}
\end{frame}

\begin{frame}{Results: Language Neuron Distribution for LAPE }
	\begin{block}{}\scriptsize
	  \begin{figure}[ht]
	  	\centering
	  	\begin{tabular}{cc}
	  		\includegraphics[width=0.3\textwidth]{Chapters/Chapter_5/lang_neurons/Meta-Llama-3.1-8B/lape/set1/lang_neuron_dist.png} &
	  		\includegraphics[width=0.3\textwidth]{Chapters/Chapter_5/lang_neurons/Meta-Llama-3.1-8B/lape/set3/lang_neuron_dist.png}
	  	\end{tabular}
	  \end{figure}
	\end{block}
\end{frame}

\begin{frame}{Results: Language Neuron Overlap for LAPE }
	\begin{block}{}\scriptsize
		\begin{figure}[ht]
			\centering
			\begin{tabular}{cc}
				\includegraphics[width=0.425\textwidth]{Chapters/Chapter_5/lang_neurons/Meta-Llama-3.1-8B/lape/set1/lang_neurons_overlap.png} &
				\includegraphics[width=0.425\textwidth]{Chapters/Chapter_5/lang_neurons/Meta-Llama-3.1-8B/lape/set3/lang_neurons_overlap.png}
			\end{tabular}
		\end{figure}
	\end{block}
\end{frame}

\begin{frame}{Results: Layer-wise Neuron Distribution for LAPE }
	\begin{block}{}\scriptsize
		\begin{figure}[ht]
			\centering
			\begin{tabular}{cc}
				\includegraphics[width=0.3\textwidth]{Chapters/Chapter_5/lang_neurons/Meta-Llama-3.1-8B/lape/set1/layerwise_lang_neurons_dist.png} &
				\includegraphics[width=0.3\textwidth]{Chapters/Chapter_5/lang_neurons/Meta-Llama-3.1-8B/lape/set3/layerwise_lang_neurons_dist.png}
			\end{tabular}
		\end{figure}
	\end{block}
\end{frame}

\begin{frame}{Results: Perplexity Change for LAPE }
	\begin{block}{}\scriptsize
		\begin{figure}[ht]
			\centering
			\begin{tabular}{cc}
				\includegraphics[width=0.425\textwidth]{Chapters/Chapter_5/perplexity/Meta-Llama-3.1-8B/set1/ppx_change.png} &
				\includegraphics[width=0.425\textwidth]{Chapters/Chapter_5/perplexity/Meta-Llama-3.1-8B/set3/ppx_change.png}
			\end{tabular}
		\end{figure}
	\end{block}
\end{frame}
	
\begin{frame}{Results: Language Neuron Overlap for Activation Based Method}
	\begin{block}{}\scriptsize
		\begin{figure}[ht]
			\centering
			\begin{tabular}{ccc}
				\includegraphics[width=0.425\textwidth]{Chapters/Chapter_5/lang_neurons/Meta-Llama-3.1-8B/act_prob_90p/lang_neurons_overlap.png} &
			\end{tabular}
		\end{figure}
	\end{block}
\end{frame}

\begin{frame}{Results: Layer-wise Dist \& Perplexity for Activation Based Method}
	\begin{columns}
		\column{0.25\textwidth}
		\begin{block}{}\scriptsize
			\begin{figure}[ht]
				\centering
				\begin{tabular}{c}
					\includegraphics[width=0.7\textwidth]{Chapters/Chapter_5/lang_neurons/Meta-Llama-3.1-8B/act_prob_90p/layerwise_lang_neurons_dist.png} 
				\end{tabular}
			\end{figure}
		\end{block}
		\column{0.65\textwidth}
		\begin{block}{}\scriptsize
			\begin{figure}[ht]
				\centering
				\begin{tabular}{c}
					\includegraphics[width=0.6\textwidth]{Chapters/Chapter_5/perplexity/Meta-Llama-3.1-8B/act_prob_95p/ppx_change.png}
					\end{tabular}
			\end{figure}
		\end{block}
	\end{columns}
\end{frame}

\begin{frame}{Results: Fine-Tuning of Baseline for LAPE}
	\begin{block}{\scriptsize XNLI Accuracy Results}\scriptsize
		\begin{table}[ht]
			\centering
			\renewcommand{\arraystretch}{1.3}
			\setlength{\tabcolsep}{8pt}
			\begin{tabularx}{\textwidth}{|X|X|X|X|}
				\hline
				\textbf{Baseline Type} & \textbf{Train Lang} & \textbf{Eval Lang} & \textbf{No Int Accuracy} \\
				\hline
				Zero Shot & en & vi & 0.8053 \\
				& en & hi & 0.7501 \\
				& en & ur & 0.7002 \\
				\hline
				Direct Fine-tuning & vi & vi & 0.8511 \\
				& hi & hi & 0.8043 \\
				& ur & ur & 0.6697 \\
				\hline
			\end{tabularx}
		\end{table}
	\end{block}
	\begin{block}{\scriptsize Analysis}\scriptsize
		\begin{itemize}
			\item The observation that direct fine-tuning for Urdu (ur) yields lower accuracy compared to zero-shot transfer because  the dataset used for fine-tuning might be of low quality.
			\item The pretrained model’s under-representation of Urdu during training might also limit the efficacy of fine-tuning.
			\item In contrast, zero-shot transfer utilizes cross-lingual embeddings and the model’s generalized multilingual capabilities, enabling better performance.
		\end{itemize}
	\end{block}
\end{frame}
	
\begin{frame}{Results: Targeted Neuron Fine-Tuning for LAPE}
	\begin{block}{\scriptsize XNLI Accuracy Results for Intervention by Higher Quantiles}\scriptsize
		\begin{table}[ht]
			\centering
			\renewcommand{\arraystretch}{1.3}
			\setlength{\tabcolsep}{4pt}
			\begin{tabularx}{\textwidth}{|X|X|X|X|X|X|X|}
				\hline
				\textbf{FT Neurons} & \textbf{Eval Lang} & \textbf{No Int} & \textbf{Int by $\mu$} & \textbf{Int by P75} & \textbf{Int by P90} & \textbf{Int by P95} \\
				\hline
				null & vi & 0.8053 & 0.7955 & 0.7915 & 0.7903 & 0.7897 \\
				set1\_en & vi & 0.8017 & 0.7957 & 0.7867 & 0.7847 & 0.7847 \\
				set1\_vi & vi & 0.8007 & 0.7951 & 0.7881 & 0.7855 & 0.7845 \\
				set1\_en+vi & vi & 0.8009 & 0.7943 & 0.7865 & 0.7845 & 0.7833 \\
				\hline
				null & hi & 0.7501 & 0.7521 & 0.7496 & 0.7499 & 0.7506 \\
				set1\_en & hi & 0.7476 & 0.7456 & 0.7454 & 0.7448 & 0.7434 \\
				set6\_en & hi & 0.7491 & 0.7456 & 0.7464 & 0.7458 & 0.7441 \\
				set6\_hi & hi & 0.7492 & 0.7456 & 0.7461 & 0.7448 & 0.7431 \\
				set6\_en+hi & hi & 0.7488 & 0.7448 & 0.7451 & 0.7446 & 0.7431 \\
				\hline
				null & ur & 0.7002 & 0.7044 & 0.7028 & 0.6938 & 0.6904 \\
				set6\_en & ur & 0.6984 & 0.7041 & 0.7026 & 0.6964 & 0.6923 \\
				set6\_ur & ur & 0.6984 & 0.7054 & 0.7024 & 0.6952 & 0.6913 \\
				set6\_en+ur & ur & 0.6992 & 0.7061 & 0.7021 & 0.6954 & 0.6919 \\
				\hline
			\end{tabularx}
		\end{table}
	\end{block}
\end{frame}
	
\begin{frame}{Results: Targeted Neuron Fine-Tuning for LAPE}
	\begin{block}{\scriptsize XNLI Accuracy Results for Intervention by Lower Quantiles}\scriptsize
		\begin{table}[ht]
			\centering
			\renewcommand{\arraystretch}{1.3}
			\setlength{\tabcolsep}{4pt}
			\begin{tabularx}{\textwidth}{|X|X|X|X|X|X|X|X|X|}
				\hline
				\textbf{FT Neurons} & \textbf{Eval Lang} & \textbf{No Int} & \textbf{Int by P5} & \textbf{Int by P10} & \textbf{Int by P25} & \textbf{Int by P50} \\
				\hline
				null & vi & 0.8053 & 0.7366 & 0.7773 & 0.8055 & 0.7971 \\
				set1\_en & vi & 0.8017 & 0.7432 & 0.7803 & 0.8027 & 0.7923 \\
				set1\_vi & vi & 0.8007 & 0.7426 & 0.7801 & 0.8029 & 0.7927 \\
				set1\_en+vi & vi & 0.8009 & 0.7426 & 0.7805 & 0.8027 & 0.7919 \\
				\hline
				null & hi & 0.7501 & 0.7488 & 0.7518 & 0.7496 & 0.7526 \\
				set1\_en & hi & 0.7476 & 0.7451 & 0.7468 & 0.7461 & 0.7451 \\
				set6\_en & hi & 0.7491 & 0.7464 & 0.7456 & 0.7476 & 0.7468 \\
				set6\_hi & hi & 0.7492 & 0.7461 & 0.7456 & 0.7464 & 0.7456 \\
				set6\_en+hi & hi & 0.7488 & 0.7456 & 0.7468 & 0.7462 & 0.7454 \\
				\hline
				null & ur & 0.7002 & 0.6861 & 0.6873 & 0.6971 & 0.7054 \\
				set6\_en & ur & 0.6984 & 0.6886 & 0.6904 & 0.7016 & 0.7031 \\
				set6\_ur & ur & 0.6984 & 0.6873 & 0.6913 & 0.7011 & 0.7034 \\
				set6\_en+ur & ur & 0.6992 & 0.6867 & 0.6893 & 0.7016 & 0.7042 \\
				\hline
			\end{tabularx}
		\end{table}
	\end{block}
\end{frame}
	
\begin{frame}{Results: Targeted Neuron Fine-Tuning for Act\_Prob\_90p}
	\begin{block}{\scriptsize XNLI Accuracy Results for Intervention by Higher Quantiles}\scriptsize
		\begin{table}[ht]
			\centering
			\renewcommand{\arraystretch}{1.3}
			\setlength{\tabcolsep}{4pt}
			\begin{tabularx}{\textwidth}{|X|X|X|X|X|X|X|}
				\hline
				\textbf{FT Neurons} & \textbf{Eval Lang} & \textbf{No Int} & \textbf{Int by $\mu$} & \textbf{Int by P75} & \textbf{Int by P90} & \textbf{Int by P95} \\
				\hline
				null & vi & 0.8053 & 0.7821 & 0.7891 & 0.7931 & 0.7957 \\
				en & vi & 0.8027 & 0.7801 & 0.7865 & 0.7907 & 0.7923 \\
				vi & vi & 0.8005 & 0.7799 & 0.7857 & 0.7901 & 0.7907 \\
				en+vi & vi & 0.8005 & 0.7799 & 0.7853 & 0.7897 & 0.7909 \\
				\hline
				null & hi & 0.7501 & 0.7411 & 0.7268 & 0.7186 & 0.7122 \\
				en & hi & 0.7492 & 0.7386 & 0.7262 & 0.7214 & 0.7148 \\
				hi & hi & 0.7486 & 0.7382 & 0.7262 & 0.7212 & 0.7148 \\
				en+hi & hi & 0.7488 & 0.7372 & 0.7256 & 0.7214 & 0.7144 \\
				\hline
				null & ur & 0.7002 & 0.6971 & 0.6978 & 0.6932 & 0.6875 \\
				en & ur & 0.6994 & 0.7006 & 0.6978 & 0.6937 & 0.6893 \\
				ur & ur & 0.6988 & 0.7001 & 0.6982 & 0.6935 & 0.6925 \\
				en+ur & ur & 0.6978 & 0.6992 & 0.6961 & 0.6927 & 0.6915 \\
				\hline
			\end{tabularx}
		\end{table}
	\end{block}
\end{frame}	
	
\begin{frame}{Results: Targeted Neuron Fine-Tuning for Act\_Prob\_90p}
	\begin{block}{\scriptsize XNLI Accuracy Results for Intervention by Lower Quantiles}\scriptsize
		\begin{table}[ht]
			\centering
			\renewcommand{\arraystretch}{1.3}
			\setlength{\tabcolsep}{4pt}
			\begin{tabularx}{\textwidth}{|X|X|X|X|X|X|X|}
				\hline
				\textbf{FT Neurons} & \textbf{Eval Lang} & \textbf{Int by P5} & \textbf{Int by P10} & \textbf{Int by P25} & \textbf{Int by P50} \\
				\hline
				null & vi & 0.7653 & 0.7747 & 0.7801 & 0.7835 \\
				en & vi & 0.7667 & 0.7749 & 0.7789 & 0.7827 \\
				vi & vi & 0.7653 & 0.7743 & 0.7779 & 0.7817 \\
				en+vi & vi & 0.7663 & 0.7745 & 0.7795 & 0.7805 \\
				\hline
				null & hi & 0.7501 & 0.7468 & 0.7422 & 0.7378 \\
				en & hi & 0.7481 & 0.7451 & 0.7426 & 0.7378 \\
				hi & hi & 0.7486 & 0.7456 & 0.7416 & 0.7374 \\
				en+hi & hi & 0.7491 & 0.7462 & 0.7428 & 0.7364 \\
				\hline
				null & ur & 0.6948 & 0.6958 & 0.6954 & 0.6966 \\
				en & ur & 0.6948 & 0.6946 & 0.6956 & 0.7001 \\
				ur & ur & 0.6946 & 0.6952 & 0.6952 & 0.6998 \\
				en+ur & ur & 0.6942 & 0.6964 & 0.6948 & 0.6998 \\
				\hline
			\end{tabularx}
		\end{table}
	\end{block}
\end{frame}	
	
\begin{frame}{Results: Targeted XNLI Task Neuron Fine-Tuning for Act\_Prob\_90p}
	\begin{block}{\scriptsize XNLI Accuracy Results for Intervention by Higher Quantiles}\scriptsize
	\begin{table}[ht]
		\centering
		\renewcommand{\arraystretch}{1.3}
		\setlength{\tabcolsep}{4pt}
		\begin{tabularx}{\textwidth}{|X|X|X|X|X|X|X|}
			\hline
			\textbf{FT Neurons} & \textbf{Eval Lang} & \textbf{No Int} & \textbf{Int by $\mu$} & \textbf{Int by P75} & \textbf{Int by P90} & \textbf{Int by P95} \\
			\hline
			null & vi & 0.8053 & 0.7926 & 0.7994 & 0.8014 & 0.8028 \\
			en & vi & 0.7956 & 0.7926 & 0.8004 & 0.8006 & 0.8026 \\
			vi & vi & 0.7951 & 0.7926 & 0.7998 & 0.7998 & 0.8014 \\
			en+vi & vi & 0.7921 & 0.7924 & 0.7982 & 0.7921 & 0.8012 \\
			\hline
		\end{tabularx}
	\end{table}
	\end{block}
	\begin{block}{\scriptsize XNLI Accuracy Results for Intervention by Lower Quantiles}\scriptsize
		
		\begin{table}[ht]
			\centering
			\renewcommand{\arraystretch}{1.3}
			\setlength{\tabcolsep}{4pt}
			\begin{tabularx}{\textwidth}{|X|X|X|X|X|X|X|}
				\hline
				\textbf{FT Neurons} & \textbf{Eval Lang} & \textbf{Int by P5} & \textbf{Int by P10} & \textbf{Int by P25} & \textbf{Int by P50} \\
				\hline
				null & vi & 0.7706 & 0.7801 & 0.7912 & 0.7936 \\
				en & vi & 0.7702 & 0.7791 & 0.7896 & 0.7926 \\
				vi & vi & 0.7688 & 0.7774 & 0.7896 & 0.7931 \\
				\hline
			\end{tabularx}
		\end{table}
	\end{block}
\end{frame}	


\begin{frame}{Conclusions and Future Directions}
	\begin{columns}
		\column{0.45\textwidth}
		\begin{block}{\scriptsize Conclusions}\scriptsize
			\begin{itemize}
				\item Statistical activation patterns (e.g., mean, P75, P90, P95) yielded limited consistent improvements across configurations.
				\item Language-specific neurons derived from Wikipedia data provide strong baselines for zero-shot transfer.
				\item Zero-shot transfer performance occasionally surpassed direct fine-tuning, particularly in languages like Urdu.
				\item Interventions using lower quantiles (e.g., P5, P10) degraded performance, while higher quantiles were more effective.
			\end{itemize}		  
		\end{block}
		\column{0.45\textwidth}
		\begin{block}{\scriptsize Future Directions}\scriptsize
			\begin{itemize}
				\item Investigate neuron selection at the attention layer for improved fine-tuning strategies.
				\item Analyze cross-lingual generalization impacts, especially for low-resource languages.
				\item Scale methods to larger architectures like Meta-Llama-13B for insights into scalability.
				\item Explore interpretability to understand why zero-shot transfer sometimes outperforms direct fine-tuning.
			\end{itemize}
		\end{block}
	\end{columns}
\end{frame}


	\section{References}
	\begin{frame}{References}
		\begin{block}{}\scriptsize
			\printbibliography
		\end{block}
	\end{frame}
	
	\begin{frame}{}
		\centering \Huge
		\emph{Thank You!}
	\end{frame}
	
\end{document}